\section{Introduction}
\label{sec:intro}

Coding agents work through ordinary file-system tools: they traverse a
repository, inspect source files, edit code, and run builds and
tests~\cite{sweagent2024,openhands2025}. Consider a service that launches
several attempts from the same repository revision and dependency set.
Each attempt needs a consistent starting view and must be able to change
files independently. Making that starting state available is only the first
part of serving a task. Throughout execution, the file system must also
resolve names and attributes, locate file contents, and serve reads and writes.

Sharing a baseline addresses provisioning, but leaves these access costs
open. Layered namespaces and copy-on-write branches already let tasks reuse
state while keeping changes private~\cite{overlayfs,branchfs2026,deltafs2021}.
A shared directory tree still needs a metadata access path, and its files
still need a data access path. Immutable distribution systems support
indexed, demand-loaded state~\cite{cvmfs2011,nydus_rafs}; grouping records
can amortize requests, but may also fetch records a task never uses. On the
data path, larger objects can amortize object requests, while smaller
independently readable units can limit excess transfer and verification.
The problem is therefore to support independent task views while controlling
the work of provisioning state, obtaining metadata, and accessing file data.

Our observations motivate considering access throughout a task's lifetime.
We reanalyzed 1,138,884 FUSE requests from eight coding-agent task runs
(Section~\ref{sec:agent-observations}). Metadata operations, including
unsuccessful probes, made up more than half of the requests in each of the
two workload groups. Reads and writes had mixed request lengths, and
activity varied substantially across tasks and over time. A provisioning
mechanism alone leaves this metadata and data traffic to the underlying
storage path.

Our design separates the state a task shares from the units used to access
that state. A workspace needs an immutable baseline identity and rules for
resolving private changes; those rules need not prescribe the baseline's
metadata format. Immutable metadata can be organized for lookup and reuse
whether or not its consumer needs a writable workspace. Likewise, object
formation and backend read units can be selected without changing which
file version a task sees. These boundaries allow choices driven by state
reuse, metadata locality, and data access to remain distinct.

\sys realizes this design in an object-backed file system with three
mechanisms. \emph{Delta-Based Workspaces} resolve a shared immutable
baseline together with private namespace and data-mapping changes, avoiding
a separate materialized baseline for each task. \emph{Packed Metadata}
places immutable names, attributes, and data locators in indexed, bounded
groups on object storage so that readers can fetch and reuse related
records. \emph{Variable-Granularity Object Storage} lets deployments choose
client object size separately from optional encoding and verification
layouts on a compatible backend, to trade request count against transferred
bytes and processing work.

The mechanisms are optional and can be composed. A workspace can inherit
a native metadata layer; a packed snapshot can serve a read-only consumer;
and an ordinary writable volume can choose its object block size with the
backend's default layout. When tasks reuse a stable environment, multiple
workspaces can instead reference one packed baseline. This composition
shares the baseline's representation while preserving each task's private
changes; it does not make packing a prerequisite for workspace isolation.

The paper develops the state and access rules that make these choices
compatible: private-view resolution, immutable metadata lookup and
publication, and the mapping from file ranges to object and backend work
units. Section~\ref{sec:implementation} identifies the paths present in the
current code and their integration limits. The completed workload study
characterizes request structure; it does not establish cross-task sharing
or system speedups. Performance and failure experiments remain to be
completed under the evaluation plan in Section~\ref{sec:evaluation}.
